\section{The \texorpdfstring{$\{3,4,3,4\}$}{3,4,3,4} Honeycomb}

With the fourth dimension and the $24$-cell in hand, the whole substrate can now be described in plain language. The seemingly preferred substrate, the one this paper argues for, is the $\{3,4,3,4\}$ hyperbolic honeycomb. The selection itself is made later, on grounds rather than by assumption, so here the geometry is fixed and laid out, not yet justified. This section sets the stage the rest of the paper lives on, the four layers of bulk, coin, cusp, and boundary.

\subsection{In essence}

The honeycomb $\{3,4,3,4\}$ fills 4D hyperbolic space completely with identical $24$-cells, in the way cubes fill ordinary 3D space. It is regular, meaning every cell, face, edge, and corner is alike, and it is hyperbolic, meaning it fits only in negatively curved 4-space. Its most important feature is that its corners sit infinitely far away. They are \emph{ideal} corners, and the pattern you see looking out toward one of them is the ordinary cubic grid of flat 3D space. So $\{3,4,3,4\}$ is a curved 4D crystal whose edge, seen up close, is flat three-dimensional space.

\subsection{How to read the symbol}

Each number in the symbol adds a dimension by recording how the previous shape is assembled.

\begin{itemize}
  \item $\{3\}$ is a triangle.
  \item $\{3,4\}$ says four triangles around each corner close into an octahedron.
  \item $\{3,4,3\}$ says three octahedra around each edge close into the $24$-cell.
  \item $\{3,4,3,4\}$ says four $24$-cells around each triangular face fill 4D hyperbolic space, the honeycomb.
\end{itemize}

Two facts read straight off the symbol. The cell, found by dropping the last number, is $\{3,4,3\}$, the $24$-cell. The vertex figure, found by dropping the first number, is $\{4,3,4\}$, the cubic honeycomb. The second of these is the magic. The corner of $\{3,4,3,4\}$ is ordinary 3D cubic space.

\subsection{The four layers}

The substrate is best held in mind as four layers, each carrying a distinct role.

The \emph{bulk} is 4D hyperbolic space tiled by $\{3,4,3,4\}$. Being hyperbolic it grows exponentially with radius, each shell about ten times the last, and that growth will be the cosmic expansion. It is paracompact rather than compact, which the cusp below turns into a feature.

The \emph{coin} is the set of $24$ neighbour directions at each cell. They form the $24$-cell, whose vertices are the $D_4$ root system. This is the single most important structural fact in the paper. The $24$ directions decompose under the symmetry as $8v + 8s + 8c$, the $\mathrm{SO}(8)$ vector plus its two spinors, with the $S_3$ triality permuting the three eights. The coin is where spin, the gauge group, and the three generations come from.

The \emph{cusp} is the vertex figure $\{4,3,4\}$, the cubic honeycomb, which is Euclidean, infinite, and flat. This is the paracompact ideal element, and it is physical space, a flat 3D cubic cusp where matter, light, and gravity propagate. So physical space is three-dimensional and flat, and it is the cusp of the 4D curved bulk. The bulk is the fire, the cusp is the radiance it casts.

The \emph{boundary} is the $3$-sphere $S^3$ at infinity, the holographic screen. Gravity and quantum nonlocality live in the relation between bulk and boundary. Time enters as the beat together with the radial direction, which doubles as cosmic time and as the renormalization axis.

\subsection{The vertex figure is flat 3D space}

This is the whole point, so it is worth stating slowly. In a compact honeycomb the vertex figure is a finite shape and the corner is an ordinary point. But $\{4,3,4\}$ is infinite, since it is all of flat 3D space. When the vertex figure is infinite, the corner is pushed out to infinity, an ideal vertex. So the corners of $\{3,4,3,4\}$ live infinitely far away, and the cross-section around one of those ideal corners is exactly the cubic grid of flat 3D Euclidean space.

\begin{result}[The cusp is exact flat 3D space]
Because the vertex figure of $\{3,4,3,4\}$ is the infinite cubic honeycomb $\{4,3,4\}$, every vertex is ideal and the cross-section at each ideal vertex is the flat 3D cubic lattice. Physical space is recovered with no approximation, as the exact vertex figure rather than as an extracted slice.
\end{result}

\subsection{Why paracompactness is a feature}

A natural worry is that paracompactness, the presence of corners at infinity, is some kind of defect. It is the opposite. A regular honeycomb gives a clean flat physical layer only if it has a parabolic, that is Euclidean, element to foliate by. A compact honeycomb, such as $\{5,3,4\}$ or $\{7,3\}$, has none. A horosphere cutting through it crosses its cells aperiodically, producing a thin disordered slab that is rough and sub-three-dimensional. This is why the compact substrates have no clean flat space. The ideal $\{4,3,4\}$ vertex figure of $\{3,4,3,4\}$ instead gives a clean periodic cubic cross-section. So paracompactness is precisely what supplies the orderly 3D space.

In the standard taxonomy of hyperbolic Coxeter tilings, sorted by the signature of the Gram matrix of mirror angles, the classes run Euclidean, compact or Lannér, paracompact or Koszul, hypercompact or Vinberg, and Lorentzian. Paracompact means having at least one ideal vertex. The honeycomb $\{3,4,3,4\}$ is paracompact, and the crystallographic spinor-bearing honeycombs are always in the paracompact or flat classes, never the compact one \cite{margenstern2013}.

\subsection{Hyperbolic, and built by reflection}

The honeycomb does not fit in flat 4-space. The angles do not add up there. It closes only in negatively curved $H^4$, a fact decided by the Gram-matrix signature of its mirror angles. Being hyperbolic, it grows exponentially. Count the $24$-cells within radius $r$ and the number explodes like $e^{cr}$, each shell larger than everything inside it, which is the cosmic expansion. Almost all of the honeycomb lives out near its $3$-sphere boundary at infinity, and the structure there, around the ideal corners, is the cubic 3D grid.

It is generated the same way as any regular honeycomb, by reflection. Start with one fundamental chamber and reflect it across its mirror walls forever. The reflections form the Coxeter group $[3,4,3,4]$, which has five mirrors and the diagram $\circ\,3\,\circ\,4\,\circ\,3\,\circ\,4\,\circ$. Infinite complexity grows from a tiny local rule. A deep general principle sits at the boundary. The Euclidean cubic lattice at infinity is the affine, Euclidean-motion symmetry of the hyperbolic bulk. Flat space appears at infinity as the asymptotic symmetry of curved space.

\subsection{How to picture it}

No one directly visualizes a 4D hyperbolic honeycomb, so you climb a ladder of analogies. One dimension down, in flat 3D, cubes tile space as $\{4,3,4\}$ and the corner is just a point. Go one dimension up and bend the space negatively, and $\{3,4,3,4\}$ is the analogue whose corner is a whole copy of that cubic tiling. The Escher disk gives a second image. It is like the hyperbolic disk where figures crowd toward the rim, but in 4D, with $24$-cells crowding toward a $3$-sphere boundary, and the crowding pattern at the rim spelling out cubic 3D space. The cusp funnel gives a third. Look down an infinitely long funnel in the honeycomb, and far down it the cross-section settles into a fixed cubic grid that never changes shape and only recedes. That steady cubic cross-section is the flat 3D space at the ideal corner. What you can hold in mind is $24$-cell bricks, curved 4D mortar, ideal corners at infinity, and a cubic-grid cross-section at each corner.

\subsection{The control and the substrate}

Throughout the paper $\{5,3,4\}$, the dodecahedral honeycomb of 3D hyperbolic space, serves as a control. It is
the closest three-dimensional analogue of the substrate, close enough to picture (Figure~\ref{fig:dodecagrid}),
and comparing the two makes the case for $\{3,4,3,4\}$ concrete.

\begin{table}[ht]
\centering
\begin{tabular}{@{}lll@{}}
\toprule
feature & $\{5,3,4\}$, the 3D control & $\{3,4,3,4\}$, the preferred substrate \\
\midrule
lives in & hyperbolic 3-space & hyperbolic 4-space \\
cell & dodecahedron $\{5,3\}$ & $24$-cell $\{3,4,3\}$ \\
neighbours per cell & 12 & 24 \\
vertex figure & octahedron, finite & cubic honeycomb $\{4,3,4\}$, infinite \\
corners & ordinary points & ideal, at infinity \\
flat slice & flat 2D & flat 3D, the cubic grid \\
boundary at infinity & $2$-sphere & $3$-sphere \\
spinor coin & no, icosahedral $1+3+3'+5$ & yes, $D_4$ as $8v+8s+8c$ \\
\bottomrule
\end{tabular}
\caption{The 3D control honeycomb against the 4D preferred substrate.}
\end{table}

The headline is that $\{5,3,4\}$ gives flat 2D slices while $\{3,4,3,4\}$ gives flat 3D slices. That one extra dimension is exactly what ordinary physical space needs. And only the $24$-cell coin carries the spinor. The icosahedral coin of $\{5,3,4\}$ decomposes as $1 + 3 + 3' + 5$, which carries no spinor, while the $D_4$ coin decomposes as $8v + 8s + 8c$, which does \cite{pollard2026vibetest}.


This is the cleanest candidate for where flat 3D physical space comes from. It is the cubic cusp of a 4D hyperbolic bulk, recovered with no approximation, paired with a spinor-carrying cell. Which findings of the paper live on which layer is the subject of the later robustness audit. Spin and the gauge structure live in the bulk coin, matter propagates on the cusp, and the two are connected by a projection developed in the matter chapters.
